The OSPI interface must be configured to QuadSPI to be compatible with the external QSPI memory:
\begin{table}[H]
    \centering
    \begin{tabular}{ll}
        \hline
        \textbf{Parameter} & \textbf{Value} \\
        \hline
        Mode & Quad SPI \\
        Clock & Port 1 CLK \\
        Device Size & Port1 NCS \\
        Data [3:0] & Port 1 IO[3:0] \\
        Fifo Threshold & 1 \\
        Dual Quad Mode & Disable \\
        Memory Type & Micron \\
        Device Size & 24 \\
        Chip Select High Time & 3 \\
        Free Running Clock & Disable \\
        Clock Mode & Low \\
        Wrap Size & Not Supported \\
        Clock Prescaler & 5 \\
        Sample Shifting & Half Cycle \\
        Delay Hold Quarter Cycle & Disable \\
        Chip Select Boundary & 0 \\
        Maximum Transfer & 0 \\
        Refresh Rate & 0 \\
        Status (DLYB) & Bypassed \\
        \hline
    \end{tabular}
    \caption{QSPI generic configuration parameters from STM32CubeIDE.}
\end{table}
Navigate to the 'Clock Configuration' menu and locate the PLL2 section. Then, set the following:
\begin{table}[H]
    \centering
    \begin{tabular}{ll}
        \hline
        \textbf{Parameter} & \textbf{Value} \\
        \hline
        PLL3 Source Mux & HSE \\
        PLL3M & / 10 \\
        *N & X 133 \\
        /R & /2 \\
        /Q & /4 \\
        \hline
    \end{tabular}
    \caption{QSPI Clock Configuration.}
\end{table}
You should now have $266$ MHz on 'PLL2R' and $133$ MHz on PLL2Q (the QSPI clock).